\subsection{The quadratic cost of an edge-density deficit}\label{sec:nonexceptional-quadratic-growth}

At the phase boundary, we show that the reduced objective
$I_{\pc(r),r}(\eps)$ increases quadratically from the constant-graphon
cost as the edge density $\eps$ decreases from $r$, while the $H$-density
remains fixed at $r^m$.  We express this in terms of the edge-density deficit
\[
  \delta:=r-\eps.
\]
Note that by \Cref{lem:edge-density-deficit}, $\delta>0$ for every optimizer on the
symmetry-breaking side.
Fix $r_0\in(0,1)\setminus\{r_*\}$ and let $I$ be the open interval containing $r_0$
supplied by \Cref{lem:fixed-density-bipodality}, with compact closure in
$(0,1)\setminus\{r_*\}$.  We may shrink $I$ around $r_0$ below; all estimates
are uniform for $r\in I$.
At the boundary value $p=\pc(r)$, define the boundary excess
\begin{equation}\label{eq:boundary-excess}
  G_r(\delta)=I_{\pc(r),r}(r-\delta)-J_{\pc(r)}(r)
  \qquad (0<\delta<\rho_0),
\end{equation}
with $\rho_0$ as in \Cref{lem:fixed-density-bipodality}; that lemma ensures
that $I_{\pc(r),r}(r-\delta)$ is finite throughout this range.  The value
at $\delta=0$ belongs instead to the analytic extension $G(r,\delta)$
constructed in \Cref{thm:positive-second-variation}; it is not part of the domain of $G_r$.
Thus, $G_r(\delta)$ measures this additional cost relative to the constant
graphon.  More precisely, we will prove that
\[
  G_r(\delta)=\frac12A_H(r)\delta^2+O_{H,I}(\delta^3)
  \qquad (\delta\downarrow0),
  \qquad A_H(r)>0.
\]

For $r\in I$, recall that $\sm(r)$ is the second contact density associated
with the boundary point $(\pc(r),r)$.  Let $\ell_r$ be the supporting line to
the graph of $x\mapsto J_{\pc(r)}(x^{1/d})$ at $x=r^d$:
\[
  \ell_r(x)=J_{\pc(r)}(r)
  +\frac{J_{\pc(r)}'(r)}{d r^{d-1}}(x-r^d)
  \qquad (0 \le x \le 1).
\]
Define the \emph{supporting cost gap} by
\begin{equation}\label{eq:supporting-cost-gap}
  g_r(z)=J_{\pc(r)}(z)-\ell_r(z^d)
  \qquad (0\le z\le1).
\end{equation}
Since $\ell_r$ is a supporting line, $g_r\ge0$.  We claim that there is
$\gamma_{d,I}>0$ such that, uniformly for $r\in I$,
\begin{equation}\label{eq:supporting-gap-quadratic-bound}
  g_r(z)\ge \gamma_{d,I}\,\dist(z,\{r,\sm(r)\})^2.
\end{equation}
Apply \eqref{eq:contact-quadratic-separation} on $\overline I$.
The contact densities are uniformly bounded below by some $\eta_{d,I}>0$,
and for $z\in[0,1]$ and $v\in\{r,\sm(r)\}$,
\[
  |z^d-v^d|
  =
  |z-v|\,(z^{d-1}+z^{d-2}v+\cdots+v^{d-1})
  \ge
  \eta_{d,I}^{d-1}|z-v|.
\]
Taking the minimum over the two contact densities proves
\eqref{eq:supporting-gap-quadratic-bound}.

\begin{lemma}[Scalar quadratic lower bound]\label{lem:scalar-quadratic-bound}
There exist constants $C_{d,I}>0$ and $\delta_0>0$ such that, for every
$r\in I$, the following holds.  If $0<\delta<\delta_0$ and $X$ is a random variable with
values in $[0,1]$ satisfying
$\EE X=r-\delta$ and
$\EE X^d\ge r^d$,
then
\begin{equation}\label{eq:scalar-quadratic-bound}
  \EE J_{\pc(r)}(X)\ge J_{\pc(r)}(r)+C_{d,I}\delta^2.
\end{equation}
\end{lemma}

\begin{proof}
The slope of $\ell_r$ is positive, because $\pc(r)<r$.  Hence
\[
  \EE\ell_r(X^d)
  =   \ell_r(\EE X^d)
  \ge \ell_r(r^d)
  =   J_{\pc(r)}(r).
\]
Using \eqref{eq:supporting-cost-gap},
\[
  \EE g_r(X)
  =   \bigl(\EE J_{\pc(r)}(X)-\EE \ell_r(X^d)\bigr)
  \le \bigl(\EE J_{\pc(r)}(X)-J_{\pc(r)}(r)\bigr).
\]
It remains to show $\EE g_r(X)\ge C_{d,I}\delta^2$.
Set $s:=\sm(r)$ and
\[
  k:=\frac{s^d-r^d}{s-r}
  =s^{d-1}+s^{d-2}r+\cdots+r^{d-1}
  \ge\eta_{d,I}^{d-1}>0.
\]
Let $Y$ be the measurable nearest-point projection of $X$ onto $\{r,s\}$,
with ties assigned to $r$.  Since $Y\in\{r,s\}$, we have
$Y^d-r^d=k(Y-r)$.  Using $\EE X=r-\delta$ and the fact that
$z\mapsto z^d$ is $d$-Lipschitz on $[0,1]$, we obtain
\begin{align*}
  0\le\EE X^d-r^d
  &\le\EE Y^d-r^d+d\,\EE|X-Y|\\
  &=k\bigl(\EE(Y-X)-\delta\bigr)+d\,\EE|X-Y|\\
  &\le(k+d)\EE|X-Y|-k\delta.
\end{align*}
Hence
\[
  \EE|X-Y|\ge\frac{k}{k+d}\delta
  \ge\frac{\eta_{d,I}^{d-1}}{d+\eta_{d,I}^{d-1}}\delta.
\]
Since $|X-Y|=\dist(X,\{r,s\})$, the bound
\eqref{eq:supporting-gap-quadratic-bound} and Cauchy--Schwarz give
\[
  \EE g_r(X)\ge\gamma_{d,I}\,\EE|X-Y|^2
  \ge\gamma_{d,I}\bigl(\EE|X-Y|\bigr)^2
  \ge C_{d,I}\delta^2,
\]
where $C_{d,I}:=\gamma_{d,I}\bigl(\eta_{d,I}^{d-1}/(d+\eta_{d,I}^{d-1})\bigr)^2>0$.
\end{proof}

\begin{corollary}[Quadratic lower bound for graphons]\label{prop:graphon-quadratic-bound}
For every $r\in I$ and $0<\delta<\delta_0$, if $e(W)=r-\delta$ and
$t(H,W)\ge r^m$, then
\begin{equation}\label{eq:graphon-quadratic-bound}
  I_{\pc(r)}(W)
  \ge
  J_{\pc(r)}(r)+C_{d,I}\delta^2,
\end{equation}
with $C_{d,I}$ and $\delta_0$ as in \Cref{lem:scalar-quadratic-bound}.
\end{corollary}

\begin{proof}
Let $X=W(U,V)$, where $(U,V)$ is uniformly distributed on $[0,1]^2$.  Then
$\EE X=e(W)=r-\delta$.  By the generalized H\"older inequality
\eqref{eq:generalized-holder} and feasibility,
\[
  \EE X^d=\int W^d\ge t(H,W)^{d/m} \ge r^d.
\]
By \Cref{lem:scalar-quadratic-bound}, we conclude the proof.
\end{proof}

We now prove that the preceding quadratic lower bound has
the correct order.  For small $\delta>0$, we construct a bipodal graphon with
edge density $r-\delta$, $H$-density $r^m$, and $I_{\pc(r)}$-value only
$O_{H,I}(\delta^2)$ above that of the constant graphon.  The construction starts
from the constant graphon $r$, makes a small block, uses the second contact density
$s=\sm(r)$ on the cross edges, and then adjusts the large-block density to
restore the two constraints.

\begin{lemma}[Quadratic upper bound]\label{lem:bipodal-quadratic-bound}
After shrinking $I$ around $r_0$, there exist constants $C_{H,I}>0$ and
$\delta_0>0$ such that, for every $r\in I$ and $0<\delta<\delta_0$, there is a bipodal graphon
$V_\delta$ satisfying $e(V_\delta)=r-\delta$, $t(H,V_\delta)=r^m$,
and
\begin{equation}\label{eq:bipodal-quadratic-bound}
  I_{\pc(r)}(V_\delta)
  \le
  J_{\pc(r)}(r)+C_{H,I}\delta^2.
\end{equation}
\end{lemma}

\begin{proof}
Fix $a\in(0,1)$; we take $a=1/2$ at the end.
For $r$ near $r_0$, put $s=\sm(r)$ and let $V_{r,c,q}$ have block size $c$
and edge densities $a,s,q$ on the small block, cross edges, and large block,
respectively.  To solve the two constraints, define
\[
  F(r,\delta,c,q)
  :=\bigl(e(V_{r,c,q})-(r-\delta),\,t(H,V_{r,c,q})-r^m\bigr).
\]
The density formulas are polynomial in $c,a,s,q$, so they define $F$ also
for real $c$ near $0$.  Since $\sm$ is analytic near $r_0$, this extension
is analytic and satisfies $F(r,0,0,r)=0$.

Writing $n=|V(H)|=2m/d$, the assignments with zero or one vertex in the
small block give
\begin{align*}
  e(V_{r,c,q})&=c^2a+2c(1-c)s+(1-c)^2q,\\
  t(H,V_{r,c,q})&=(1-c)^nq^m
  +nc(1-c)^{n-1}s^dq^{m-d}+O_H(c^2).
\end{align*}
Thus, with
\[
  D_d(r,s):=\frac{s^d-r^d}{d r^{d-1}}-(s-r)>0,
\]
where positivity follows from strict convexity and $s\ne r$, the Jacobian is
\begin{gather*}
  \partial_{(c,q)}F(r,0,0,r)=
  \begin{pmatrix}
    2(s-r) & 1 \\
    \frac{2m}{d}r^{m-d}(s^d-r^d) & m r^{m-1}
  \end{pmatrix},\\
  \det\partial_{(c,q)}F(r,0,0,r)=-2m r^{m-1}D_d(r,s)<0.
\end{gather*}
The analytic implicit function theorem at $(r_0,0,0,r_0)$ therefore gives
jointly analytic functions $c(r,\delta)$ and $q(r,\delta)$ near $(r_0,0)$.
Local uniqueness and $F(r,0,0,r)=0$ give $c(r,0)=0$ and $q(r,0)=r$, while
\begin{equation}\label{eq:comparison-graphon-constraints}
  e(V_{r,c(r,\delta),q(r,\delta)})=r-\delta,
  \qquad
  t(H,V_{r,c(r,\delta),q(r,\delta)})=r^m.
\end{equation}
Shrink $I$ and choose $\delta_0>0$ so that
$\overline I\times[-\delta_0,\delta_0]$ lies in this analytic domain.
Differentiating the constraints at $\delta=0$ and using the Jacobian gives,
with uniform Taylor remainders,
\begin{equation}\label{eq:comparison-block-size}
  c(r,\delta)=\frac{\delta}{2D_d(r,s)}+O_{H,I,a}(\delta^2),
\end{equation}
\begin{equation}\label{eq:comparison-block-density}
  q(r,\delta)=r-\frac{2(s^d-r^d)}{d r^{d-1}}c(r,\delta)
  +O_{H,I,a}(c(r,\delta)^2).
\end{equation}
Since $D_d(r,\sm(r))$ is bounded above and away from zero on $\overline I$,
we may decrease $\delta_0$ so that $0<c(r,\delta)<1/2$ and $0<q(r,\delta)<1$
for every $r\in I$ and $0<\delta<\delta_0$.
Hence $V_\delta:=V_{r,c(r,\delta),q(r,\delta)}$ is a graphon with the required
densities.

It remains to estimate its cost.  Abbreviate
$c=c(r,\delta)$, $q=q(r,\delta)$.  Since $q-r=O_{H,I,a}(c)$ and all edge
densities and $\pc(r)$ stay away from $0$ and $1$, Taylor's theorem gives
\begin{align*}
  I_{\pc(r)}(V_\delta)-J_{\pc(r)}(r)
  &=J_{\pc(r)}'(r)(q-r)+2c\{J_{\pc(r)}(s)-J_{\pc(r)}(r)\}
    +O_{H,I,a}(c^2)\\
  &=2c\left\{J_{\pc(r)}(s)-J_{\pc(r)}(r)
    -\frac{J_{\pc(r)}'(r)}{d r^{d-1}}(s^d-r^d)\right\}
    +O_{H,I,a}(c^2).
\end{align*}
The bracket vanishes because $\ell_r$ touches the scalar curve at both
$r^d$ and $s^d$; see (M3)--(M4) of \Cref{thm:scalar-lz-boundary}.
Thus the cost difference is $O_{H,I,a}(c^2)=O_{H,I,a}(\delta^2)$ by
\eqref{eq:comparison-block-size}.  Taking $a=1/2$ proves
\eqref{eq:bipodal-quadratic-bound}.
\end{proof}

Recall that the boundary excess \eqref{eq:boundary-excess} is defined only for
$\delta>0$, because the Kenyon--Radin--Ren--Sadun regime of \Cref{thm:krrs-analytic-extension}
requires a strictly positive $H$-density excess $\tau-\eps^m$.  The theorem
below therefore produces, in addition to the quadratic expansion, the
real-analytic \emph{extension} of $G_r$ to a two-sided neighborhood of
$\delta=0$; it is that extension whose $\delta$-derivatives appear in
\Cref{sec:nonexceptional-proof}.

\begin{theorem}[Universal positivity of the scalar second variation]\label{thm:positive-second-variation}
After shrinking $I$ around $r_0$, there exist $\overline\delta>0$, an open
set $\mathcal N$ containing $I\times(-\overline\delta,\overline\delta)$,
and a real-analytic
function $G:\mathcal N\to\RR$ such that
\begin{equation}\label{eq:boundary-excess-extension}
  G(r,\delta)=G_r(\delta)
  \qquad(r\in I,\ 0<\delta<\overline\delta).
\end{equation}
Setting $A_H(r):=\partial_\delta^2G(r,0)$, the following hold.
\begin{enumerate}[label=(\roman*)]
  \item $G(r,0)=0$ and $\partial_\delta G(r,0)=0$ for every $r\in I$.
  \item $A_H$ is real-analytic on $I$ and satisfies $A_H(r)\ge a_0>0$
  there for some constant $a_0$.
  \item There is $C_2\ge1$, depending only on $H,I$, such that
  \begin{equation}\label{eq:boundary-excess-curvature}
    \bigl|\partial_\delta^2G(r,\delta)-A_H(r)\bigr|\le C_2|\delta|
    \qquad(r\in I,\ |\delta|<\overline\delta).
  \end{equation}
  \item Consequently the boundary excess has the expansion
  \begin{equation}\label{eq:boundary-excess-expansion}
    G_r(\delta)=\frac12A_H(r)\delta^2+O_{H,I}(\delta^3)
    \qquad(r\in I,\ 0<\delta<\overline\delta).
  \end{equation}
\end{enumerate}
\end{theorem}

\begin{proof}
Use the single analytic parameter map from
\Cref{thm:krrs-analytic-extension} chosen in the proof of
\Cref{lem:fixed-density-bipodality}.  Evaluating it at
\[
  (\eps,\vartheta)=(r-\delta,\,r^m-(r-\delta)^m)
\]
gives jointly analytic parameters $(q_{11},q_{12},q_{22},c)$ near $(r_0,0)$,
including small negative $\delta$.  Define
\[
  G(r,\delta):=
  c^2J_{\pc(r)}(q_{11})+2c(1-c)J_{\pc(r)}(q_{12})
  +(1-c)^2J_{\pc(r)}(q_{22})-J_{\pc(r)}(r).
\]
The three edge-density parameters and $\pc(r)$ stay in $(0,1)$ near
$(r_0,0)$, so $G$ is analytic there.  For $\delta>0$, this is the cost of
the fixed-density optimizer minus the constant-graphon cost, by
\eqref{eq:reduced-objective-attainment}.  At $\delta=0$, the parameters have
$c=0$ and $q_{22}=r$, giving $G(r,0)=0$.
Shrink $I$ and choose $\overline\delta>0$ so that
$\overline I\times[-\overline\delta,\overline\delta]$ lies in the analytic
domain and the preceding upper and lower bounds apply throughout the
positive half of this rectangle.

For $r\in I$ and $0<\delta<\overline\delta$, the fixed-density optimizer
satisfies \Cref{prop:graphon-quadratic-bound} and has no greater cost than
the graphon in \Cref{lem:bipodal-quadratic-bound}.  Consequently,
\[
  C_{d,I}\delta^2\le G_r(\delta)\le C_{H,I}\delta^2.
\]
Together with analyticity and $G(r,0)=0$, this gives
$\partial_\delta G(r,0)=0$ and
\[
  A_H(r)=\partial_\delta^2G(r,0)
  =\lim_{\delta\downarrow0}\frac{2G_r(\delta)}{\delta^2}>0.
\]
Thus $A_H$ is positive and analytic near $r_0$; after shrinking $I$,
continuity gives $A_H(r)\ge a_0:=A_H(r_0)/2>0$ for all $r\in I$.

The continuous function $|\partial_\delta^3G|$ is bounded on the closed
rectangle above.  Taking $C_2\ge1$ larger than this bound, the mean value
theorem gives \eqref{eq:boundary-excess-curvature}.  Taylor's theorem then
gives \eqref{eq:boundary-excess-expansion}, with remainder bounded by
$C_2|\delta|^3/6$.  All bounds are uniform for $r\in I$.
\end{proof}

The function $A_H$ does not depend on the chosen neighborhood or analytic
extension: any two extensions agree for small positive $\delta$, where
both equal \eqref{eq:boundary-excess}, and hence their second derivatives
at $\delta=0$ agree.  Since $r_0$ was arbitrary, these local functions
define a positive real-analytic function
\begin{equation}\label{eq:second-variation-domain}
  A_H:(0,1)\setminus\{r_*\}\longrightarrow(0,\infty).
\end{equation}
